\documentclass[12pt,a4paper]{article}
\usepackage[margin=25mm, headheight=15pt, headsep=10pt]{geometry}
\usepackage{fontspec}
\usepackage[table]{xcolor} % Note: table option is required for rowcolors
\usepackage{titlesec}
\usepackage{tocloft}
\usepackage{fancyhdr}
\usepackage{booktabs}
\usepackage{tcolorbox}
\tcbuselibrary{skins, breakable}
\usepackage[colorlinks=true, linkcolor=emerald, urlcolor=emerald, bookmarks=true]{hyperref}

\definecolor{emerald}{RGB}{0,102,51}
\definecolor{darkred}{RGB}{139,0,0}
\definecolor{lightgray}{RGB}{245,245,245}

\usepackage[nil]{babel}
\babelprovide[import, main]{bengali}
\babelprovide[import]{arabic}
\babelfont{rm}[Renderer=HarfBuzz,Scale=1.0]{Noto Serif Bengali}
\babelfont{sf}[Renderer=HarfBuzz]{Noto Sans Bengali}
\babelfont{tt}[Renderer=HarfBuzz]{Noto Sans Bengali}
\babelfont[arabic]{rm}[Renderer=HarfBuzz,Scale=1.1]{Noto Naskh Arabic}

% Section styling
\titleformat{\section}{\Large\bfseries\color{emerald}}{\thesection.}{0.5em}{}
\titleformat{\subsection}{\large\bfseries\color{emerald!85!black}}{\thesubsection.}{0.5em}{}
\titleformat{\subsubsection}{\normalsize\bfseries\color{emerald!70!black}}{\thesubsubsection.}{0.5em}{}
\titlespacing*{\section}{0pt}{18pt}{8pt}

% Fancy Header/Footer setup
\pagestyle{fancy}
\fancyhf{} % clear all header and footer fields
\fancyhead[L]{\color{emerald}\small\textbf{\nouppercase{\leftmark}}}
\fancyfoot[L]{\scriptsize\color{black!50}{\lat{AI}}-সংকলিত, আলেম-যাচাইকৃত নয়}
\fancyfoot[C]{\color{emerald!70!black}\thepage}
\renewcommand{\headrulewidth}{0.4pt}
\renewcommand{\headrule}{\hbox to\headwidth{\color{emerald}\leaders\hrule height \headrulewidth\hfill}}
\renewcommand{\footrulewidth}{0pt}

% Custom boxes
% 1. Ayat/Hadith Box
\newtcolorbox{ayatbox}[1][]{
    enhanced, breakable, 
    colback=emerald!5, 
    colframe=emerald, 
    boxrule=0.5pt, 
    arc=3pt, 
    left=10pt, right=10pt, top=10pt, bottom=10pt,
    #1
}

% 2. Quote/Translation Box
\newtcolorbox{quotebox}[1][]{
    enhanced, breakable,
    frame hidden,
    colback=lightgray,
    borderline west={3pt}{0pt}{emerald},
    left=15pt, right=10pt, top=10pt, bottom=10pt,
    sharp corners,
    #1
}

% 3. AI-generated notice box
\newtcolorbox{warnbox}[1][]{
    enhanced, breakable,
    colback=red!4,
    colframe=darkred,
    boxrule=0.8pt,
    arc=3pt,
    left=10pt, right=10pt, top=10pt, bottom=10pt,
    #1
}

% Commands
\newcommand{\arab}[1]{\foreignlanguage{arabic}{#1}}
\newcommand{\kib}[1]{\textcolor{emerald}{\textbf{#1}}}

% Latin (English) fallback: Noto Serif Bengali has NO Latin glyphs
\newfontfamily\latfont[Renderer=HarfBuzz]{Noto Serif}
\newcommand{\lat}[1]{{\latfont #1}}

\setlength{\parskip}{5pt}
\setlength{\parindent}{0pt}

% Fix for missing bullet in Noto Serif Bengali
\renewcommand{\labelitemi}{$\bullet$}



\begin{document}
\begin{center}{\Large\bfseries\color{emerald} \lat{03}-আয়াতুল-কুরসি-এবং-অন্যান্য}\end{center}
\vspace{1em}
\section*{আয়াতুল কুরসি ও অন্যান্য গুরুত্বপূর্ণ আয়াত — শব্দ-শব্দ অর্থ}

\begin{warnbox}
 \textbf{গুরুত্বপূর্ণ নোটিশ (\lat{Important} \lat{Notice}):} এই কন্টেন্টটি \lat{AI} (কৃত্রিম বুদ্ধিমত্তা) দিয়ে প্রস্তুত ও সংকলিত। \lat{AI} কঠোর সূত্র-মান নীতি মেনে শুধু নির্ভরযোগ্য উৎস থেকে তথ্য সংগ্রহ করেছে — কেবল সহীহ/হাসান হাদিস ও সহীহ সনদের আসার রাখা হয়েছে, দুর্বল সনদ বাদ দেওয়া হয়েছে। তবে এটি কোনো আলেম বা মুহাদ্দিস কর্তৃক যাচাইকৃত নয়।
\end{warnbox}

\medskip\hrule\medskip

\subsection*{আয়াত ১: আয়াতুল কুরসি (সূরা আল-বাকারা ২:২৫৫)}

\subsubsection*{মূল আরবি}
\begin{center}\arab{اللَّهُ لَا إِلَٰهَ إِلَّا هُوَ الْحَيُّ الْقَيُّومُ ۚ لَا تَأْخُذُهُ سِنَةٌ وَلَا نَوْمٌ ۚ لَهُ مَا فِي السَّمَاوَاتِ وَمَا فِي الْأَرْضِ ۗ مَنْ ذَا الَّذِي يَشْفَعُ عِنْدَهُ إِلَّا بِإِذْنِهِ ۚ يَعْلَمُ مَا بَيْنَ أَيْدِيهِمْ وَمَا خَلْفَهُمْ ۖ وَلَا يُحِيطُونَ بِشَيْءٍ مِنْ عِلْمِهِ إِلَّا بِمَا شَاءَ ۚ وَسِعَ كُرْسِيُّهُ السَّمَاوَاتِ وَالْأَرْضَ ۖ وَلَا يَئُودُهُ حِفْظُهُمَا ۚ وَهُوَ الْعَلِيُّ الْعَظِيمُ}\end{center}

\subsubsection*{বাংলা উচ্চারণ}
\textbf{আল্লাহু লা ইলাহা ইল্লা হুয়াল হাইয়ুল কাইয়ুম। লা তা'খুযুহু সিনাতুন ওয়ালা নাউম। লাহু মা ফিস সামাওয়াতি ওয়া মা ফিল আরয। মান যাল্লাজী ইয়াশফাউ ইন্দাহু ইল্লা বিইযনিহি। ইয়া'লামু মা বাইনা আইদীহিম ওয়া মা খালফাহুম। ওয়ালা ইউহীতুনা বিশাই'মিন ইলমিহি ইল্লা বিমা শায়া। ওয়াসিআ কুরসিয়্যুহুস সামাওয়াতি ওয়াল আরয। ওয়ালা ইয়া'উদুহু হিফযুহুমা। ওয়া হুয়াল আলিয়্যুল আযীম।}

\subsubsection*{শব্দ-শব্দ অর্থ}

\begin{table}[h]\centering\small
\begin{tabular}{lll}
আরবি & বাংলা & \lat{English} \\
\hline
\textbf{\arab{اللَّهُ}} & আল্লাহ & \textit{\lat{Allah}} \\
\textbf{\arab{لَا} \arab{إِلَٰهَ}} & কোনো ইলাহ/উপাস্য নেই & \textit{\lat{there} \lat{is} \lat{no} \lat{god}} \\
\textbf{\arab{إِلَّا} \arab{هُوَ}} & কিন্তু তিনি & \textit{\lat{except} \lat{Him}} \\
\textbf{\arab{الْحَيُّ}} & চিরঞ্জীব & \textit{\lat{the} \lat{Ever-Living}} \\
\textbf{\arab{الْقَيُّوم}ু} & চিরস্থায়ী, সমস্ত সৃষ্টি তাঁর উপর নির্ভরশীল & \textit{\lat{the} \lat{Sustainer} \lat{of} \lat{all}} \\
\textbf{\arab{لَا} \arab{تَأْخُذُهُ}} & তাঁকে গ্রাস করে না & \textit{\lat{Neither} \lat{overtakes}} \\
\textbf{\arab{سِنَةٌ}} & ঘুম/মূর্ছা & \textit{\lat{drowsiness}} \\
\textbf{\arab{وَلَا} \arab{نَوْمٌ}} & ও ঘুমও নয় & \textit{\lat{nor} \lat{sleep}} \\
\textbf{\arab{لَهُ}} & তাঁরই & \textit{\lat{For} \lat{Him}} \\
\textbf{\arab{مَا} \arab{فِي}} & যা আছে & \textit{\lat{is} \lat{whatever}} \\
\textbf{\arab{السَّمَاوَاتِ}} & আকাশে/পরলোকে & \textit{\lat{in} \lat{the} \lat{heavens}} \\
\textbf{\arab{وَمَا} \arab{فِي}} & ও যা আছে & \textit{\lat{and} \lat{whatever}} \\
\textbf{\arab{الْأَرْضِ}} & পৃথিবীতে & \textit{\lat{in} \lat{the} \lat{earth}} \\
\textbf{\arab{مَنْ} \arab{ذَا}} & কে & \textit{\lat{Who} \lat{is} \lat{it}} \\
\textbf{\arab{يَشْفَعُ}} & সুপারিশ করে & \textit{\lat{that} \lat{intercedes}} \\
\textbf{\arab{عِنْدَهُ}} & তাঁর কাছে & \textit{\lat{before} \lat{Him}} \\
\textbf{\arab{إِلَّا} \arab{بِإِذْنِهِ}} & তাঁর অনুমতি ছাড়া & \textit{\lat{except} \lat{by} \lat{His} \lat{permission}} \\
\textbf{\arab{يَعْلَمُ}} & তিনি জানেন & \textit{\lat{He} \lat{knows}} \\
\textbf{\arab{مَا} \arab{بَيْنَ} \arab{أَيْدِيهِمْ}} & তাদের সামনের বিষয় & \textit{\lat{what} \lat{is} \lat{before} \lat{them}} \\
\textbf{\arab{وَمَا} \arab{خَلْفَهُمْ}} & ও পেছনের বিষয় & \textit{\lat{and} \lat{what} \lat{is} \lat{behind} \lat{them}} \\
\textbf{\arab{وَلَا} \arab{يُحِيطُونَ}} & তারা জানে না/গ্রাস করতে পারে না & \textit{\lat{They} \lat{encompass}} \\
\textbf{\arab{بِشَيْءٍ}} & কোনো কিছু & \textit{\lat{anything}} \\
\textbf{\arab{مِنْ} \arab{عِلْمِهِ}} & তাঁর জ্ঞান থেকে & \textit{\lat{of} \lat{His} \lat{knowledge}} \\
\textbf{\arab{إِلَّا} \arab{بِمَا} \arab{شَاءَ}} & তিনি যা চান তা ছাড়া & \textit{\lat{except} \lat{what} \lat{He} \lat{wills}} \\
\textbf{\arab{وَسِعَ}} & বিস্তৃত/পরিপূর্ণ & \textit{\lat{Encompasses}} \\
\textbf{\arab{كُرْسِيُّهُ}} & তাঁর কুরসী (সিংহাসন/কর্তৃত্ব) & \textit{\lat{His} \lat{Throne} (\lat{authority})} \\
\textbf{\arab{السَّمَاوَاتِ}} & আকাশসমূহকে & \textit{\lat{the} \lat{heavens}} \\
\textbf{\arab{وَالْأَرْضَ}} & ও পৃথিবীকে & \textit{\lat{and} \lat{the} \lat{earth}} \\
\textbf{\arab{وَلَا} \arab{يَئُودُهُ}} & তাঁকে কষ্ট দেয় না & \textit{\lat{Wearies} \lat{Him} \lat{not}} \\
\textbf{\arab{حِفْظُهُمَا}} & তাদের সংরক্ষণ & \textit{\lat{the} \lat{preservation} \lat{of} \lat{both}} \\
\textbf{\arab{وَهُوَ}} & এবং তিনই & \textit{\lat{And} \lat{He} \lat{is}} \\
\textbf{\arab{الْعَلِيُّ}} & সর্বোচ্চ & \textit{\lat{the} \lat{Most} \lat{High}} \\
\textbf{\arab{الْعَظِيمُ}} & মহান & \textit{\lat{the} \lat{Most} \lat{Great}} \\
\hline
\end{tabular}
\end{table}

\subsubsection*{পূর্ণ বাংলা অনুবাদ}
\textbf{\lat{“}আল্লাহ — তিনি ব্যতীত কোনো ইলাহ নেই; তিনি চিরঞ্জীব, চিরস্থায়ী (সমস্ত সৃষ্টি তাঁর উপর নির্ভরশীল)। তাঁকে ঘুমও গ্রাস করে না, মূর্ছাও নয়। আকাশ ও পৃথিবীতে যা কিছু আছে সব তাঁরই। কে তাঁর কাছে সুপারিশ করতে পারে তাঁর অনুমতি ছাড়া? তিনি জানেন তাদের সামনের ও পেছনের বিষয়। তারা তাঁর জ্ঞানের কোনো কিছুই গ্রাস করতে পারে না — তিনি যা চান তা ছাড়া। তাঁর কুরসী (কর্তৃত্ব) আকাশ ও পৃথিবীকে বিস্তৃত। তাদের সংরক্ষণ তাঁকে ক্লান্ত করে না। আর তিনি সর্বোচ্চ, মহান।\lat{”}}

\subsubsection*{সংক্ষিপ্ত ব্যাখ্যা}
আয়াতুল কুরসি কুরআনের \textbf{সবচেয়ে গুরুত্বপূর্ণ আয়াত} হিসেবে বিবেচিত। হাদিসে এসেছে: \textbf{\lat{“}নামাজের পর আয়াতুল কুরসি পড়লে জান্নাতে প্রবেশের আদেশ দেওয়া হয়।\lat{”}} (সুনানে তিরমিযী ৩৫২০, সুনানে আবু দাউদ ৫০১৫ — সহীহ)। এটি \textbf{\lat{tawheed} (তাওহীদ)}, \textbf{\lat{power} (ক্ষমতা)}, \textbf{\lat{knowledge} (জ্ঞান)} ও \textbf{\lat{mercy} (রহমাত)}-এর পূর্ণাঙ্গ ঘোষণা।

\medskip\hrule\medskip

\subsection*{আয়াত ২: আল-কুরআনুল কারিম (সূরা কাহফ ১৮:১০৩-১১০ — শেষ দশ আয়াত)}

\subsubsection*{শেষ ৩টি আয়াত (১৮:১০৮-১১০)}

\begin{center}\arab{إِنَّمَا الْأَعْمَالُ بِالنِّيَّاتِ}\end{center}

\begin{quote}
\textbf{\lat{“}নিশ্চয়ই সকল আমল নিয়তের ওপর নির্ভরশীল।\lat{”}} — সহীহ বুখারি ১, সহীহ মুসলিম ১৯০৭ (আলেই সিলসিলা হাদিস)
\end{quote}

এটি \textbf{হাদিসে কুদসি} নয়, বরং মুহাম্মদ \arab{ﷺ}-এর বাণী (মারফু')। কিন্তু নামাজে প্রবেশের আগে \textbf{\lat{intention} (নিয়ত)}-র গুরুত্ব স্মরণ করার জন্য এটি অপরিহার্য।

\medskip\hrule\medskip

\subsection*{সংক্ষিপ্ত ব্যাখ্যা ও ব্যবহার}

\subsubsection*{আয়াতুল কুরসি — কখন পড়বেন?}

\begin{table}[h]\centering\small
\begin{tabular}{ll}
সময় & ফজিলত \\
\hline
\textbf{ফরজ নামাজের পর} & সুন্নাতে মুআক্কাদাহ — প্রতিটি ফরজ নামাজের পর ১ বার \\
\textbf{ইশার ফরজের পর} & আওয়াবীন নামাজের আগে — উত্তম সময় \\
\textbf{তাহাজ্জুদের পর} & বিশেষ মাকাম \\
\textbf{ঘুমানোর আগে} & সুরক্ষা ও আল্লাহর স্মরণ \\
\textbf{বাড়ি থেকে বের হওয়ার আগে} & সুরক্ষা \\
\hline
\end{tabular}
\end{table}

\subsubsection*{শেষ দশ আয়াত — কখন পড়বেন?}

\begin{table}[h]\centering\small
\begin{tabular}{ll}
সময় & ফজিলত \\
\hline
\textbf{শেষ দুই রাকাত (ইশার সান্ধ্য)} & সুন্নাতে মুআক্কাদাহ \\
\textbf{সকাল-বিকালের ২ রাকাত হায়াতুল উম্মাহ} & সুন্নাতে মুআক্কাদাহ \\
\textbf{প্রতিটি ২ রাকাত নফলে} & সুন্নাত \\
\hline
\end{tabular}
\end{table}

\subsubsection*{প্রয়োগের পরামর্শ}

\begin{enumerate}
\item \textbf{মুখস্থ করুন} — আয়াতুল কুরসি ও শেষ দশ আয়াত।
\item \textbf{অর্থ মনে করুন} — নামাজে পড়ার সময়।
\item \textbf{নিয়ত জোরালো করুন} — প্রতিটি কাজের আগে।
\item \textbf{সুরক্ষার জন্য ব্যবহার করুন} — ভয়, উদ্বেগ, বা অসুস্থতায়।
\item \textbf{নামাজের পর পড়ুন} — জান্নাতের প্রতিশ্রুতি অর্জনের জন্য।
\end{enumerate}
\medskip\hrule\medskip

\textit{শেষ আপডেট: আগস্ট ২০২৬}

\end{document}
